% ===============================================
% ECEV 33400: Stochastic Processes      Fall 2017
% hw_draft.tex
% Template for first draft of homework submission
% ===============================================


% -----------------------------------------------
% The preamble that follows can be ignored. Go on
% down to the section that says "START HERE" 
% -----------------------------------------------

\documentclass{article}

\usepackage[margin=1in]{geometry} 
\usepackage{amsmath,amsthm,amssymb}
\usepackage[sharp]{easylist}

\newcommand{\R}{\mathbb{R}}
\newcommand{\Z}{\mathbb{Z}}
\newcommand{\N}{\mathbb{N}}
\newcommand{\Q}{\mathbb{Q}}
\newcommand{\C}{\mathbb{C}}
\newcommand{\F}{{\mathbb F}}
\newcommand{\Pb}{{\mathbb P}}
\newcommand{\E}{{\mathbb E}}
\newcommand{\Frac}[2]{\left(\frac{#1}{#2}\right)}
\newcommand{\Eval}[2]{\big|_{#1}^{#2}}
\newcommand{\intinf}{\int_{-\infty}^{\infty}}

\newenvironment{theorem}[2][Theorem]{\begin{trivlist}
\item[\hskip \labelsep {\bfseries #1}\hskip \labelsep {\bfseries #2.}]}{\end{trivlist}}
\newenvironment{lemma}[2][Lemma]{\begin{trivlist}
\item[\hskip \labelsep {\bfseries #1}\hskip \labelsep {\bfseries #2.}]}{\end{trivlist}}
\newenvironment{exercise}[2][Exercise]{\begin{trivlist}
\item[\hskip \labelsep {\bfseries #1}\hskip \labelsep {\bfseries #2.}]}{\end{trivlist}}
\newenvironment{problem}[2][Problem]{\begin{trivlist}
\item[\hskip \labelsep {\bfseries #1}\hskip \labelsep {\bfseries #2.}]}{\end{trivlist}}
\newenvironment{question}[2][Question]{\begin{trivlist}
\item[\hskip \labelsep {\bfseries #1}\hskip \labelsep {\bfseries #2.}]}{\end{trivlist}}
\newenvironment{corollary}[2][Corollary]{\begin{trivlist}
\item[\hskip \labelsep {\bfseries #1}\hskip \labelsep {\bfseries #2.}]}{\end{trivlist}}

\newenvironment{solution}{\begin{proof}[Solution]}{\end{proof}}


\begin{document}

% ------------------------------------------ %
%                 START HERE                 %
% ------------------------------------------ %

\title{Homework 4} % Replace X with the appropriate number
\author{ \iffalse STUDENT NAME \fi \\ECEV 33400} % Replace  "\iffalse X \fi" with your name.
\date{Due Date: Thursday 11/02/2017} % Don't change this! :( 

\maketitle

% -----------------------------------------------------
% The following two environments (theorem, proof) are
% where you will enter the statement and proof of your
% first problem for this assignment.
%
% In the theorem environment, you can replace the word
% "theorem" in the \begin and \end commands with
% "exercise", "problem", "lemma", etc., depending on
% what you are submitting. Replace the "x.yz" with the
% appropriate number for your problem.
%
% If your problem does not involve a formal proof, you
% can change the word "proof" in the \begin and \end
% commands with "solution".
% -----------------------------------------------------

\subsection*{Problems} 
% You can copy these problems down to the environments below to write your solutions & delete the top subsections.

\begin{easylist}[enumerate]
# Consider the dynamics of a closed population driven by density-independent birth and death rates. The growth (or decline) of such a population ca be described in terms of a one-step process with following transition rates:
\begin{eqnarray}
g_n & = & \beta\,n  \nonumber \\ 
r_n & = & \delta\,n \nonumber
\end{eqnarray}
## Write the master equation to describe the temporal evolution for the probability of the population attaining $n$ individuals at time $t$. 
## Consider the population have $n_0$ individuals at the initial time. Write the equation for the temporal evolution of the probability generating function, $G(x,t)$, and solve it using the method of characteristics. 
## Calculate the expected extinction time of the population. Consider the case that the population have $n_0$ individuals at the initial time and that the population will certainly decline, this is, $\beta < \delta$. Check through simulations that your expression matches the average of extinction time over stochastic replicates. 

# Consider the simplest stochastic model underlying Island Biogegoraphy Theory. This momdel describes the temporal evolution for the number of species in an island, $n$. Species colonize the island at rate $\lambda$ from a mainland and undergo extinction. The mainland harbors a total number of potentially colonizing species, $S$.   
## Write the master equation to describe the temporal evolution for the probability of the island harboring $n$ species at time $t$. 
## Consider the island contains no species at the initial time. Write the equation for the temporal evolution of the probability generating function, $G(x,t)$, and solve it using the method of characteristics.
## Give an explicit expression for the temporal evolution of variance in the number of species in the island.
## Usually bigger islands harbor more species. Assume that extinction an colonization rates scale with area according to:
\begin{eqnarray}
\lambda & = & \lambda_0\, A^{\beta}   \nonumber \\ 
\delta  & = & \delta_0 \, A^{-\alpha} \nonumber  
\end{eqnarray}
where $\alpha$ and $\beta$ are positive exponents. Calculate the model prediction for the species-area relationship, this is, the function giving the number of species in the island as area size increases. Check your result trough simulations.   
\end{easylist}

% -----------------------------------------------------
% First problem
% -----------------------------------------------------

% \begin{problem}{1}
% Show that the probability-generating function for the binomial distribution $p_j = \binom{n}{j} p^j q^{n-j}$, where $ q = 1 - p$, is $P(x) = (q + px)^n$. Hence prove that the mean is $m = np$, and the variance is $\sigma^2 = npq$.
% \end{problem}

% \begin{solution}
% \end{solution}

% -----------------------------------------------------
% Second problem
% -----------------------------------------------------

% \vspace{0.25in} % This just adds some space between problems 1 and 2.

% \begin{problem}{2}
% What is the probability-generating function for the Poisson distribution $$p_j = \frac{e^{-\lambda} \lambda^j}{j!}, j = 0, 1, 2, \cdots ?$$ Hence show that the mean and variance of the distribution are both equal to $\lambda$. 
% \end{problem}

% \begin{solution}
% \end{solution}

% -----------------------------------------------------
% Third problem
% -----------------------------------------------------


% -----------------------------------------------------
% Fourth problem
% -----------------------------------------------------


% -----------------------------------------------
% Ignore everything that appears below this.
% -----------------------------------------------

\end{document}
